% Generated draft: requires mathematical and source review.
\documentclass[11pt]{article}
\usepackage[margin=25mm]{geometry}
\usepackage{fontspec}
\setmainfont{Latin Modern Roman}
\usepackage{amsmath,amssymb,mathtools,unicode-math}
\setmathfont{Latin Modern Math}
\usepackage{enumitem,needspace,xurl,hyperref}
\setlength{\parindent}{0pt}
\setlength{\parskip}{5pt}
\setlength{\emergencystretch}{3em}
\newcommand{\N}{\mathbb{N}}
\newcommand{\Z}{\mathbb{Z}}
\newcommand{\Q}{\mathbb{Q}}
\newcommand{\R}{\mathbb{R}}
\newcommand{\C}{\mathbb{C}}
\newcommand{\F}{\mathbb{F}}
\newcommand{\A}{\mathbb{A}}
\newcommand{\PP}{\mathbb{P}}
\newcommand{\E}{\mathbb{E}}
\newcommand{\eps}{\varepsilon}
\newcommand{\dd}{\,\mathrm{d}}
\begin{document}
\section*{UnsolvedMath ID: 30002820}
\textbf{Source problem code:} catalogue entry
\subsection{Definitions and mathematical statement}
Let \(T\) be a triangulated surface, with vertex set \(V(T)\) and edge set \(E(T)\). Write \(\lambda_{ij}>0\) and \(\mu_{ij}>0\) for the lengths assigned by two metrics to an edge \(\{i,j\}\in E(T)\). The discrete conformal equivalence specified in the record is vertex scaling: there exists a function \(u:V(T)\to\R\) such that
\[
\mu_{ij}=\exp\!\left(\frac{u_i+u_j}{2}\right)\lambda_{ij}
\qquad\text{for every }\{i,j\}\in E(T).
\]
Here \(u_i=u(i)\). The source does not specify the geometric category or the additional admissibility conditions required of these edge lengths as metrics.

The supplied question asks whether there exists a metric subdivision scheme that preserves this equivalence. In a general preservation formulation, whenever \(T,\lambda,\mu\) satisfy the preceding hypotheses, the scheme should produce metrics \(\lambda'\) and \(\mu'\) on a common subdivision \(T'\) of \(T\), for which there exists \(u':V(T')\to\R\) satisfying
\[
\mu'_{ab}=\exp\!\left(\frac{u'_a+u'_b}{2}\right)\lambda'_{ab}
\qquad\text{for every }\{a,b\}\in E(T').
\]
This formulation records the requested preservation property, but the class of allowed subdivision schemes remains unresolved in the supplied record. In particular, it is not stated whether the same rule must be applied separately to each metric, whether subdivision must be nontrivial, how subdivided lengths must relate to the original metric, or whether \(u'\) must extend \(u\). No finiteness, compactness, or boundary hypotheses on \(T\) are supplied. Consequently, these formulas do not by themselves specify a fully determined existence problem.
\subsection{Short English statement}
The question asks for a subdivision procedure that preserves discrete conformal equivalence expressed by scaling each edge using factors at its two endpoints. The supplied record leaves the permitted metric geometry and subdivision rules unspecified, so those qualifications require clarification before the existence question has a precise scope.
\subsection{Sources}
\begin{itemize}
\item Supplied local UnsolvedMath registry record, permanent ID \(30002820\), exact source problem code \texttt{catalogue entry}, titled \emph{Discrete-Conformal Metric Subdivision Schemes}. The record places the problem in the Oberwolfach Reports Open Problems collection and gives proposed year 2015; no individual proposer is supplied.
\item Supplied source citation: \emph{Discrete Differential Geometry} (2015), Oberwolfach Reports, DOI 10.4171/owr/2015/13, pp.\ 661--729. Supplied source and literature link, labeled \emph{OWR Report 13/2015}: \url{https://doi.org/10.4171/owr/2015/13}. The background describes these reports as documenting mathematical workshops and their open-problem discussions.
\item Supplied literature source: \emph{Discrete conformal maps and ideal hyperbolic polyhedra}, \url{https://arxiv.org/abs/1005.2698}.
\item Supplied literature assessment dated exactly \texttt{2026-08-21}: the record labels the problem open and reports that no general metric subdivision scheme preserving the stated discrete conformal equivalence was found proved or disproved in the checked literature. This is supplied dated curation derived from the OWR record, not an independently verified literature assessment; the time-sensitive status and cited claims require independent verification.
\end{itemize}
\end{document}
